% Owner: Aram
% No preamble here: only the section body. Packages and colours come from ../main.tex.
% Cite sources with keys from ../refs.bib (e.g. \cite{urvlect2}).
% Guidance: summary, limitations, open questions for Task 2.
\section{Conclusions and Next Steps}\label{sec:conclusions}

\textbf{Summary.} We have designed a multi-agent system for conflict detection and resolution in a simulated airspace. The system has five crews and eight agents, coordinated by a CrewAI Flow with three routers. The main design choices are:
\begin{itemize}
    \item \emph{An independent safety crew.} The four core operational crews are complemented by a Safety Validation \& QA crew. The agent that plans a manoeuvre is therefore never the one that certifies it, and secondary conflicts created by a manoeuvre are caught before reporting.
    \item \emph{Planning matched to the environment.} The environment is accessible, deterministic and static, and it is discretised by design (Section~\ref{sec:environment}). Deliberative planning is therefore affordable, and every plan can be verified before it is reported.
    \item \emph{Tools compute, agents decide.} Every safety-relevant computation is assigned to a deterministic tool, including the two custom tools for conflict detection and manoeuvre calculation. The LLM agents coordinate the work, take the decisions and explain them.
\end{itemize}

\textbf{Limitations.} The design deliberately simplifies the problem:
\begin{itemize}
    \item \emph{Kinematics.} Aircraft fly straight legs at constant speed. Wind, climb and descent profiles are not modelled.
    \item \emph{No learning.} The agents do not use memory and do not learn between runs.
    \item \emph{No sector capacity.} The inputs contain no sector data, so the system does not check sector capacity. This is left as future work.
    \item \emph{LLM dependence.} The quality of the decisions and explanations depends on the LLM used. The design limits this risk: no manoeuvre can be reported as resolved unless it has passed the deterministic safety check.
\end{itemize}

\textbf{Open questions for Task 2.} The next phase turns this design into a CrewAI project. The main decisions still open are:
\begin{enumerate}
    \item \emph{Process per crew:} sequential or hierarchical, and which tasks can run in parallel or asynchronously.
    \item \emph{Structured outputs:} the Pydantic models exchanged between crews, starting from the intermediate files in Table~\ref{tab:crews}.
    \item \emph{Routers:} the implementation of routers A--C in the Flow, including the retry loop with $N = 2$.
    \item \emph{Test scenarios:} input files beyond the baseline sample scenarios, so that every branch of the flow is exercised. Examples are multi-aircraft crossings, a secondary conflict created by a manoeuvre, a conflict-free scenario, an invalid input and a scenario with restricted airspace.
\end{enumerate}
